% Generated draft: requires mathematical and source review.
\documentclass[11pt]{article}
\usepackage[margin=25mm]{geometry}
\usepackage{fontspec}
\setmainfont{Latin Modern Roman}
\usepackage{amsmath,amssymb,mathtools,unicode-math}
\setmathfont{Latin Modern Math}
\usepackage{enumitem,needspace,xurl,hyperref}
\setlength{\parindent}{0pt}
\setlength{\parskip}{5pt}
\setlength{\emergencystretch}{3em}
\newcommand{\N}{\mathbb{N}}
\newcommand{\Z}{\mathbb{Z}}
\newcommand{\Q}{\mathbb{Q}}
\newcommand{\R}{\mathbb{R}}
\newcommand{\C}{\mathbb{C}}
\newcommand{\F}{\mathbb{F}}
\newcommand{\A}{\mathbb{A}}
\newcommand{\PP}{\mathbb{P}}
\newcommand{\E}{\mathbb{E}}
\newcommand{\eps}{\varepsilon}
\newcommand{\dd}{\,\mathrm{d}}
\begin{document}
\section*{UnsolvedMath ID: 30004041}
\textbf{Source problem code:} catalogue entry
\subsection{Definitions and mathematical statement}
The supplied record does not specify the graph universe. The formulation below uses finite simple undirected graphs, including the graph with no vertices; this conventional interpretation requires source confirmation.

For such a graph \(G\) and a vertex \(v\in V(G)\), the local complementation \(G*v\) is obtained by replacing the subgraph induced by the neighbors of \(v\) with its complement, leaving all other adjacencies unchanged. A graph \(H\) is a vertex minor of \(G\), written \(H\leq_{\mathrm{vm}}G\), if a graph isomorphic to \(H\) can be obtained from \(G\) by a finite sequence of local complementations and vertex deletions, possibly an empty sequence.

Let \(\mathcal G\) denote this graph universe. A class \(\mathcal F\subseteq\mathcal G\) is closed under isomorphism and vertex minors if
\[
G\in\mathcal F,\quad H\cong G\ \Longrightarrow\ H\in\mathcal F,
\qquad
G\in\mathcal F,\quad H\leq_{\mathrm{vm}}G\ \Longrightarrow\ H\in\mathcal F.
\]
The class is proper if \(\mathcal F\neq\mathcal G\). No nonemptiness assumption on \(\mathcal F\) is supplied.

For this formulation, an excluded vertex minor for \(\mathcal F\) is a graph \(H\notin\mathcal F\) such that
\[
K\leq_{\mathrm{vm}}H\ \text{and}\ |V(K)|<|V(H)|
\quad\Longrightarrow\quad K\in\mathcal F.
\]
Here minimality means minimality under vertex minors with strictly fewer vertices. Local complementation alone can change the isomorphism type without changing the number of vertices, so this convention is stated explicitly and requires confirmation against the source.

\textbf{Question.} For every proper class \(\mathcal F\subsetneq\mathcal G\) closed under isomorphism and vertex minors, is the set
\[
\bigl\{[H]_{\cong}: H\text{ is an excluded vertex minor for }\mathcal F\bigr\}
\]
finite, where \([H]_{\cong}\) denotes the isomorphism class of \(H\)? Equivalently, does every such \(\mathcal F\) admit a finite, possibly empty, family \(\mathcal B\subseteq\mathcal G\) such that, for every \(G\in\mathcal G\),
\[
G\in\mathcal F
\quad\Longleftrightarrow\quad
\text{there is no }B\in\mathcal B\text{ with }B\leq_{\mathrm{vm}}G?
\]
The graph universe and excluded-minor convention above are explicit interpretive qualifications, not hypotheses verified from the cited report.
\subsection{Short English statement}
Must every proper graph class closed under isomorphism and vertex minors have only finitely many minimal forbidden graphs, up to isomorphism? The displayed formulation uses finite simple undirected graphs and minimality under vertex minors with fewer vertices; these conventions require confirmation from the source. The supplied record labels the question open.
\subsection{Sources}
\begin{itemize}
\item Local UnsolvedMath registry, \texttt{data/unsolvedmath/problems.json}, permanent ID \(30004041\), title: Finite Obstruction Sets for Vertex-Minor-Closed Graph Classes. Supplied source problem code: \texttt{catalogue entry}. Collection: Oberwolfach Reports Open Problems. No individual proposer is identified; the supplied proposed year is 2019.
\item Supplied report citation: Graph Theory (2020), Oberwolfach Reports, DOI 10.4171/owr/2019/1, pp.\ 5--63. \url{https://doi.org/10.4171/owr/2019/1}. The background describes Oberwolfach Reports as documenting mathematical workshops and their open-problem discussions. The literature-source label is Oberwolfach Report 1/2019; the supplied 2020 title date and 2019 report designation have not been reconciled.
\item Supplied literature reference: Vertex-minors survey. \url{https://sangiloum.github.io/legacy/2024/03/2023vertexminors-survey-revised.pdf}. No author or further bibliographic details are supplied for this reference.
\item Supplied literature assessment, checked on \texttt{2026-08-22}: the record labels the problem open and states that the conjectured finite forbidden-vertex-minor characterization remains open and that current surveys list it as a central unresolved problem. This is supplied dated curation derived from the supplied OWR curation, not an independently verified literature assessment. The background explicitly requests independent verification of time-sensitive status and cited claims.
\end{itemize}
\end{document}
